%% ma: L10-B22-0002
%% lop: 10
%% chuong: 7
%% bai: 22
%% dang: 10.22.2
%% loai: TLN
%% muc: VD
%% dap_an: "363"
%% nguon: "(NBV)-ĐỀ SỐ 17-MỖI NGÀY 1 ĐỀ THI 2025.pdf (D17, MỖI NGÀY 1 ĐỀ - Đề số 17), Phần III Câu 1"
%% kien_thuc: "Bài 22 (elip: tiêu điểm, bán kính qua tiêu, tiêu cự c^2 = a^2 - b^2)"
%% trang_thai: cho-duyet
%% ly_do: "Dạng mới đề xuất 10.22.2 chờ giáo viên duyệt"
%% ghi_chu: "Đáp án gốc 363 đúng (sympy: a - c ≈ 363 101 km). Bỏ ảnh trang trí (Trái Đất, Mặt Trăng)."
\begin{ex}
Mặt Trăng chuyển động quanh Trái Đất theo quỹ đạo là một elip $(E)$ mà Trái Đất là một tiêu điểm. $(E)$ có độ dài trục lớn và độ dài trục bé lần lượt là (khoảng) $768800$ km và $767619$ km. Hỏi khoảng cách ngắn nhất từ Trái Đất đến Mặt Trăng bằng bao nhiêu nghìn ki-lô-mét (làm tròn kết quả đến hàng đơn vị)?
\shortans{363}
\loigiai{Gọi $2a$, $2b$ là độ dài trục lớn, trục bé: $a=384400$, $b=383809{,}5$. Tiêu cự $c=\sqrt{a^2-b^2}\approx21298{,}5$.

Với $M(x;y)\in(E)$ và $F$ là tiêu điểm bên phải, $MF=a-\dfrac cax\ge a-c$ (vì $x\le a$), dấu bằng khi $M$ là đỉnh trên trục lớn gần $F$. Khoảng cách ngắn nhất là $a-c\approx384400-21298{,}5\approx363101$ km, tức khoảng $363$ nghìn km.}
\end{ex}
